\scanpage{49}
\mspage{24}
Pour la nécessité de b), on est ramené aussitôt au cas où $V$ est réduit à un pt, où cela se réduit à la \deleted{``conséquence conjecturale'' de la \unsure{conjecture-question} préliminaire.}

\margin{\figure{\begin{tikzpicture}\draw (0,0)--(0,-4.3);\end{tikzpicture}}}
\uline{Conséquence conjecturale} Soit $\Gamma'\subset\operatorname{Im}(E_U\longrightarrow\Gamma_{\mathbb{Q}})$ un sous-groupe ouvert, correspondant à un corps $k$ fini sur $\mathbb{Q}$, considérons un $k$-point de $U$, d'où un relèvement $\Gamma'\longrightarrow E_U$, de sorte que $\Gamma'$ opère sur $\Pi_U$. Ceci posé, on a
\[
\Pi_U^{\Gamma'^\circ}=\{1\}.
\]

On étudiera par la suite les relations entre cette \inserted{cons.} ``conjecturale'', et la précédente (d'apparence opposée !) concernant les $E_K$.

La conjecture fondamentale sur les modèles implique la conjecture fondamentale sur les corps, à condition de prendre soin, dans cette dernière, de se limiter aux hom. compatibles aux filtrations modéliques,
\margin{et il \unsure{vaut mieux} se borner à l'équivalence de a) et b) -- les conditions c) sur les centralisateurs \unsure{risquent} de passer mal à la \uline{limite}.}

Plus généralement, prenant \inserted{(\unsure{complètes} pour $U$)} des schémas, \deleted{de dim} qui sont des $\varprojlim$ de \deleted{schémas} \inserted{modèles élém.} \unsure{anabéliques|anabéliens}, avec morphismes de transition \inserted{des immersions ouvertes} affines
\margin{48}
